\chapter{Μοντελοποίηση του Κύβου}
\label{chap:model}

\section{Σκοπός του μοντέλου}
\label{sec:model-purpose}
Η κατασκευή ενός επιλυτή (\eng{solver}) για τον κύβο του Rubik δεν ξεκινά από τον αλγόριθμο αναζήτησης· προηγείται ο ορισμός του μοντέλου. Χρειάζεται πρώτα να οριστεί με ακρίβεια τι σημαίνει κατάσταση, τι σημαίνει νόμιμη κίνηση, πότε μια κατάσταση είναι φυσικά επιλύσιμη και πώς συγκρίνεται μια παραγόμενη λύση με την αρχική κατάσταση. Αν το μοντέλο είναι ασαφές, ακόμη και ένας γρήγορος επιλυτής δεν δίνει αξιόπιστο αποτέλεσμα. Μπορεί να επιστρέφει ακολουθία κινήσεων που βασίζεται σε λανθασμένη σύμβαση προσανατολισμού, να μετρά με διαφορετική μετρική ή να αποδέχεται καταστάσεις που δεν αντιστοιχούν σε πραγματικό κύβο.

Για τον λόγο αυτό η παρούσα εργασία χρησιμοποιεί ως κύρια εσωτερική αναπαράσταση το μοντέλο κυβιδίων (\eng{cubie model}), ενώ η αναπαράσταση ψηφίδων (\eng{facelets}) παραμένει διαθέσιμη για είσοδο, έξοδο και ανεξάρτητο έλεγχο. Η εννοιολογική σύγκριση των δύο αναπαραστάσεων και η αιτιολόγηση της επιλογής έχουν παρουσιαστεί στο Κεφάλαιο~\ref{chap:background}· το παρόν κεφάλαιο ορίζει τις συγκεκριμένες κωδικοποιήσεις και συμβάσεις της υλοποίησης. Οι συμβάσεις αυτές ακολουθούν την τεχνική τεκμηρίωση του Kociemba για τον αλγόριθμο δύο φάσεων~\cite{kociembaTwoPhase,kociembaImplementationDetails}.

Μια κατάσταση του 3×3×3 περιγράφεται από τέσσερις διανυσματικές ποσότητες: τη μετάθεση γωνιών \(cp\), τον προσανατολισμό γωνιών \(co\), τη μετάθεση ακμών \(ep\) και τον προσανατολισμό ακμών \(eo\). Ο προσανατολισμός των οκτώ γωνιών ορίζεται modulo 3 και των δώδεκα ακμών modulo 2. Η λυμένη κατάσταση είναι η κατάσταση όπου κάθε κυβίδιο βρίσκεται στη θέση αναφοράς του και όλοι οι προσανατολισμοί είναι μηδενικοί. Κάθε βασική κίνηση δεν είναι απλή μετάθεση θέσεων· συνδυάζει μια σταθερή μετάθεση με συγκεκριμένες μεταβολές προσανατολισμού, δηλαδή στρέψεις των γωνιών (με αριθμητική modulo 3) και αναστροφές των ακμών (modulo 2). Οι κινήσεις με απόστροφο και οι κινήσεις 180 μοιρών προκύπτουν από επαναλαμβανόμενη εφαρμογή της αντίστοιχης βασικής στροφής.

\section{Ονοματολογία και αρίθμηση κυβιδίων}
\label{sec:model-cubie-indexing}
Η σειρά των κυβιδίων στην υλοποίηση είναι σταθερή:

\begin{itemize}
  \item γωνίες: \(URF, UFL, ULB, UBR, DFR, DLF, DBL, DRB\),
  \item ακμές: \(UR, UF, UL, UB, DR, DF, DL, DB, FR, FL, BL, BR\).
\end{itemize}

Η ονοματολογία ακολουθεί την καθιερωμένη τεχνική σύμβαση του αλγορίθμου δύο φάσεων~\cite{kociembaImplementationDetails}: κάθε όνομα δηλώνει τις έδρες στις οποίες ανήκει το κυβίδιο στη λυμένη κατάσταση. Για παράδειγμα, η γωνία \(URF\) είναι η γωνία που στη λυμένη κατάσταση εφάπτεται στην άνω (\(U\)), στη δεξιά (\(R\)) και στην εμπρός (\(F\)) έδρα, ενώ η ακμή \(UF\) βρίσκεται ανάμεσα στην άνω και στην εμπρός έδρα. Οι γωνίες αριθμούνται από 0 έως 7 και οι ακμές από 0 έως 11, με τη σειρά της παραπάνω απαρίθμησης. Το Σχήμα~\ref{fig:cubie-indexing} απεικονίζει την αρίθμηση αυτή πάνω στον κύβο.

\begin{figure}[htbp]\centering
% Αρίθμηση κυβιδίων όπως ορίζεται στο src/rubik_optimal/cube.py
% (CORNER_NAMES, EDGE_NAMES):
%   γωνίες: URF, UFL, ULB, UBR, DFR, DLF, DBL, DRB (δείκτες 0-7)
%   ακμές:  UR, UF, UL, UB, DR, DF, DL, DB, FR, FL, BL, BR (δείκτες 0-11)
% Περιεχόμενο σχήματος μόνο: το περιβάλλον figure/caption ορίζεται στο κεφάλαιο.
\begin{tikzpicture}[x=1.2cm, y=1.2cm, line cap=round, line join=round]
  % Σημείωση προσανατολισμού (κοινή και για τα τρία στρώματα).
  \node[font=\scriptsize, text=black!75] at (5.4,3.6)
    {Κάτοψη των τριών στρωμάτων: πίσω έδρα (\texttt{B}) πάνω, μπροστινή (\texttt{F}) κάτω.};

  % Γωνιακά κυβίδια (8 θέσεις, δείκτες 0-7).
  \foreach \x/\y/\idx/\nm in {%
      0/2/2/ULB, 2/2/3/UBR, 0/0/1/UFL, 2/0/0/URF,
      7.8/2/6/DBL, 9.8/2/7/DRB, 7.8/0/5/DLF, 9.8/0/4/DFR}{
    \draw[draw=black!45, fill=blue!14] (\x,\y) rectangle (\x+1,\y+1);
    \node[align=center, font=\scriptsize] at (\x+0.5,\y+0.5)
      {\textbf{\idx}\\\texttt{\nm}};
  }

  % Ακμικά κυβίδια (12 θέσεις, δείκτες 0-11).
  \foreach \x/\y/\idx/\nm in {%
      1/2/3/UB, 0/1/2/UL, 2/1/0/UR, 1/0/1/UF,
      3.9/2/10/BL, 5.9/2/11/BR, 3.9/0/9/FL, 5.9/0/8/FR,
      8.8/2/7/DB, 7.8/1/6/DL, 9.8/1/4/DR, 8.8/0/5/DF}{
    \draw[draw=black!45, fill=orange!18] (\x,\y) rectangle (\x+1,\y+1);
    \node[align=center, font=\scriptsize] at (\x+0.5,\y+0.5)
      {\textbf{\idx}\\\texttt{\nm}};
  }

  % Κεντρικά κυβίδια (σταθερές θέσεις αναφοράς) και πυρήνας.
  \foreach \x/\y/\lab in {1/1/U, 4.9/2/B, 3.9/1/L, 5.9/1/R, 4.9/0/F, 8.8/1/D}{
    \draw[draw=black!45, fill=gray!16] (\x,\y) rectangle (\x+1,\y+1);
    \node[font=\scriptsize, text=black!65] at (\x+0.5,\y+0.5) {\texttt{\lab}};
  }
  \draw[draw=black!45, fill=gray!30] (4.9,1) rectangle (5.9,2);

  % Εξωτερικά περιγράμματα και τίτλοι στρωμάτων.
  \foreach \ox/\ttl in {%
      0/{Άνω στρώμα (\texttt{U})},
      3.9/{Μεσαίο στρώμα},
      7.8/{Κάτω στρώμα (\texttt{D})}}{
    \draw[black, line width=0.8pt] (\ox,0) rectangle (\ox+3,3);
    \node[font=\footnotesize] at (\ox+1.5,-0.45) {\ttl};
  }

  % Υπόμνημα χρωμάτων.
  \draw[draw=black!45, fill=blue!14] (0,-1.55) rectangle (0.3,-1.25);
  \node[anchor=west, font=\scriptsize] at (0.4,-1.4) {γωνιακά κυβίδια (0-7)};
  \draw[draw=black!45, fill=orange!18] (4.0,-1.55) rectangle (4.3,-1.25);
  \node[anchor=west, font=\scriptsize] at (4.4,-1.4) {ακμικά κυβίδια (0-11)};
  \draw[draw=black!45, fill=gray!16] (7.8,-1.55) rectangle (8.1,-1.25);
  \node[anchor=west, font=\scriptsize] at (8.2,-1.4) {κέντρα εδρών (σταθερά)};
\end{tikzpicture}

\caption{Η αρίθμηση των οκτώ γωνιακών και των δώδεκα ακμικών κυβιδίων στο μοντέλο της εργασίας, με τις θέσεις αναφοράς κάθε κυβιδίου.}
\label{fig:cubie-indexing}
\end{figure}

Η σύμβαση αυτή είναι σημαντική επειδή οι πίνακες κινήσεων και οι συντεταγμένες βασίζονται στη σταθερή σειρά των κυβιδίων. Αν αλλάξει η σειρά, αλλάζουν και οι ακέραιοι που αναπαριστούν τις συντεταγμένες. Η υλοποίηση επομένως δεν αντιμετωπίζει τα ονόματα ως απλή τεκμηρίωση· αποτελούν μέρος της σύμβασης διεπαφής που συνδέει το μοντέλο, τους κωδικοποιητές συντεταγμένων, τους επιλυτές και τους παραγόμενους πίνακες.

Η αναπαράσταση των γωνιών έχει δύο επίπεδα. Η μετάθεση \(cp\) δηλώνει ποια γωνία βρίσκεται σε κάθε θέση. Ο προσανατολισμός \(co\) δείχνει πώς είναι στραμμένη η γωνία αυτή μέσα στη θέση της. Αντίστοιχα, η μετάθεση \(ep\) και ο προσανατολισμός \(eo\) περιγράφουν τις ακμές. Οι κινήσεις αλλάζουν τις τιμές αυτές με διαφορετικό τρόπο. Στη σύμβαση της υλοποίησης, οι κινήσεις \(U\) και \(D\) μεταθέτουν κυβίδια χωρίς να μεταβάλλουν προσανατολισμούς, οι κινήσεις \(R\) και \(L\) στρέφουν γωνίες χωρίς να αναστρέφουν ακμές, ενώ οι κινήσεις \(F\) και \(B\) στρέφουν γωνίες και προκαλούν αναστροφές ακμών.

\section{Έλεγχος φυσικής εγκυρότητας}
\label{sec:model-validity}
Ο κύβος έχει τέσσερις περιορισμούς επιλυσιμότητας: πληρότητα κυβιδίων, μηδενικό άθροισμα στρέψεων γωνιών modulo 3, μηδενικό άθροισμα αναστροφών ακμών modulo 2 και συμφωνία ισοτιμίας μεταθέσεων. Οι περιορισμοί αυτοί έχουν αναλυθεί ως ιδιότητες της ομάδας του κύβου στο Κεφάλαιο~\ref{chap:background}. Στην υλοποίηση οι περιορισμοί αυτοί ελέγχονται απευθείας πάνω στα διανύσματα \(cp\), \(co\), \(ep\) και \(eo\). Ελέγχεται ότι κάθε γωνία και κάθε ακμή εμφανίζεται ακριβώς μία φορά και ότι οι προσανατολισμοί γωνιών ανήκουν στο \(\{0,1,2\}\) και των ακμών στο \(\{0,1\}\). Ελέγχεται επίσης ότι ικανοποιούνται τα δύο αθροιστικά κριτήρια και το κριτήριο ισοτιμίας.

Πέρα από τη θεωρητική τους σημασία, οι έλεγχοι αυτοί προστατεύουν τα πειράματα από μη έγκυρες εισόδους. Αν μια μη επιλύσιμη κατάσταση έφτανε στον επιλυτή, ένα αποτέλεσμα \code{timeout} (εξάντληση χρονικού ορίου) θα μπορούσε να παρερμηνευθεί ως αποτυχία του αλγορίθμου να βρει λύση, ενώ στην πραγματικότητα δεν υπάρχει λύση. Αντίστροφα, αν ο επαληθευτής (\eng{verifier}) αποδεχόταν μη φυσική αρχική κατάσταση, τα αρχεία των μετρήσεων θα έχαναν την αποδεικτική τους αξία.

Η είσοδος από ψηφίδες μετατρέπεται σε κατάσταση κυβιδίων και περνά από τον ίδιο έλεγχο. Αυτό σημαίνει ότι η διεπαφή δεν περιορίζεται στον έλεγχο του ότι η συμβολοσειρά έχει μήκος 54 και ότι τα χρώματα είναι τα αναμενόμενα. Αναγνωρίζει ποια γωνία ή ακμή αντιστοιχεί σε κάθε ομάδα ψηφίδων, υπολογίζει προσανατολισμούς και απορρίπτει περιγραφές που παραβιάζουν τους φυσικούς περιορισμούς. Η δυνατότητα αυτή υπάρχει στη διεπαφή γραμμής εντολών, η οποία δέχεται απευθείας συμβολοσειρά 54 ψηφίδων (\eng{facelet string}) ως είσοδο. Ο χρήστης μπορεί έτσι να δώσει μια πραγματική κατάσταση κύβου αντί για ακολουθία ανακατέματος (\eng{scramble}), όπως παρουσιάζεται στο Κεφάλαιο~\ref{chap:usage}.

\section{Η μετρική κινήσεων}
\label{sec:model-metric}
Η εργασία χρησιμοποιεί σε όλες τις αναζητήσεις τη μετρική μισών στροφών (\eng{half-turn metric}, HTM), στην οποία κάθε στροφή έδρας κατά \(\pm 90^\circ\) ή κατά \(180^\circ\) μετρά ως μία κίνηση, όπως ορίστηκε στο Κεφάλαιο~\ref{chap:background}. Η μετρική καταγράφεται ρητά ως HTM στα αποτελέσματα. Η συνεπής χρήση της είναι απαραίτητη, επειδή το μήκος μιας λύσης δεν έχει νόημα χωρίς κοινή σύμβαση μέτρησης του κόστους.

Η επιλογή της μετρικής καθορίζει άμεσα τον παράγοντα διακλάδωσης (\eng{branching factor}) της αναζήτησης. Το σύνολο κινήσεων περιέχει 18 σύμβολα: κάθε μία από τις έξι έδρες έχει στροφή δεξιόστροφη, αριστερόστροφη και 180 μοιρών. Στην αναζήτηση εφαρμόζονται δύο κανόνες που αποκόπτουν περιττές διαδοχές κινήσεων. Πρώτον, απαγορεύονται διαδοχικές κινήσεις στην ίδια έδρα, επειδή ακολουθίες όπως \(R\,R'\) ή \(U\,U2\) συμπτύσσονται σε μικρότερη ή ισοδύναμη ακολουθία. Δεύτερον, εφαρμόζεται κανόνας κανονικής διάταξης για τις απέναντι έδρες, οι οποίες ανήκουν στον ίδιο άξονα και των οποίων οι κινήσεις αντιμετατίθενται (η σειρά τους δεν αλλάζει το αποτέλεσμα). Για κάθε άξονα ορίζεται κανονική σειρά (\(U\) πριν από \(D\), \(R\) πριν από \(L\), \(F\) πριν από \(B\)) και απαγορεύεται κίνηση στην πρώτη έδρα αμέσως μετά από κίνηση στην απέναντί της. Οι δύο κανόνες δεν αλλάζουν τον χώρο των προσβάσιμων καταστάσεων ούτε τα βέλτιστα μήκη, επειδή από κάθε ζεύγος κινήσεων που αντιμετατίθενται αποκλείουν μόνο τη μία από τις δύο ισοδύναμες σειρές. Η ιδιότητα αυτή ισχύει για οποιοδήποτε υποσύνολο κινήσεων, άρα και για τα περιορισμένα σύνολα κινήσεων κάθε φάσης.

Η ομοιομορφία της μετρικής φαίνεται και στην παραγωγή δεδομένων. Τα ανακατέματα δημιουργούνται από την ίδια λίστα κινήσεων που χρησιμοποιούν οι επιλυτές. Ο επαληθευτής μετρά το μήκος με τον ίδιο αναλυτή κινήσεων. Τα μεταδεδομένα κάθε πίνακα καταγράφουν πόσες κινήσεις χρησιμοποιούνται για την παραγωγή του. Για τις γενικές προβολές χρησιμοποιούνται και οι 18 κινήσεις, ενώ για τις προβολές της δεύτερης φάσης του αλγορίθμου δύο φάσεων του Kociemba (βλ.~Κεφάλαιο~\ref{chap:algorithms}) χρησιμοποιείται το περιορισμένο σύνολο των 10 κινήσεων. Έτσι ο αναγνώστης μπορεί να διακρίνει αν ένας αριθμός προέρχεται από το πλήρες σύνολο HTM ή από υποσύνολο κινήσεων κάποιας φάσης.

\section{Μετατροπές μεταξύ ψηφίδων και κυβιδίων}
\label{sec:model-facelets}
Η εργασία διατηρεί και τις δύο αναπαραστάσεις, επειδή καλύπτουν διαφορετικούς σκοπούς: η μορφή ψηφίδων χρησιμεύει στην ανταλλαγή κατάστασης με τον χρήστη και με εξωτερικούς επιλυτές, ενώ η μορφή κυβιδίων εξυπηρετεί τους αλγορίθμους (βλ.~Κεφάλαιο~\ref{chap:background}). Οι αμφίδρομες μετατροπές ανάμεσα στις δύο μορφές είναι κρίσιμο στοιχείο της υλοποίησης.

Η μετατροπή από κυβίδια σε ψηφίδες χρησιμοποιείται σε κάθε γραμμή αποτελεσμάτων (\eng{benchmark row}). Το αποτέλεσμα αποθηκεύει τη συμβολοσειρά της αρχικής κατάστασης, ώστε ο επαληθευτής να μπορεί να την ανακατασκευάσει χωρίς να εμπιστευθεί τον επιλυτή που παρήγαγε τη γραμμή. Η αντίστροφη μετατροπή, από ψηφίδες σε κυβίδια, χρησιμοποιείται από τον επαληθευτή και από τη διεπαφή γραμμής εντολών. Με αυτόν τον τρόπο τα αποθηκευμένα αποτελέσματα είναι αυτοτελή: δεν χρειάζεται να κρατιέται κάποιο εσωτερικό αντικείμενο της Python μέσα στη γραμμή αποτελεσμάτων.

Η διπλή αναπαράσταση μειώνει επίσης την πιθανότητα σιωπηρού σφάλματος. Όταν μια λύση επιστρέφεται ως λίστα συμβόλων κινήσεων, ο επαληθευτής ξεκινά από την αποθηκευμένη κατάσταση ψηφίδων, τη μετατρέπει σε κυβίδια, εφαρμόζει τις κινήσεις και ελέγχει αν καταλήγει στη λυμένη κατάσταση. Αν υπάρχει ασυμφωνία στον αναλυτή, στη σύμβαση κινήσεων ή στη μέτρηση του μήκους, η επαλήθευση αποτυγχάνει.

\section{Συντεταγμένες προβολών}
\label{sec:model-coordinates}
Οι συντεταγμένες είναι ακέραιες κωδικοποιήσεις επιμέρους πτυχών (προβολών) της πλήρους κατάστασης~\cite{kociembaImplementationDetails}. Δεν αντικαθιστούν το μοντέλο κυβιδίων· χρησιμοποιούνται όταν ένας επιλυτής ή μια γεννήτρια πινάκων χρειάζεται έναν μικρό ακέραιο δείκτη πρόσβασης στον πίνακα. Η συντεταγμένη προσανατολισμού γωνιών (CO) έχει μέγεθος \(3^7=2187\), επειδή ο όγδοος προσανατολισμός καθορίζεται από το άθροισμα modulo 3. Η συντεταγμένη προσανατολισμού ακμών (EO) έχει μέγεθος \(2^{11}=2048\), επειδή η δωδέκατη αναστροφή καθορίζεται από το άθροισμα modulo 2. Η συντεταγμένη της μεσαίας φέτας (\eng{UD-slice}) έχει μέγεθος \(\binom{12}{4}=495\), επειδή καταγράφει ποιες τέσσερις ακμές ανήκουν στη μεσαία φέτα, χωρίς να λαμβάνει υπόψη τη μεταξύ τους σειρά.

Οι πλήρεις μεταθετικές συντεταγμένες γωνιών και ακμών έχουν υλοποιηθεί ως βασικά δομικά στοιχεία και δοκιμάζονται με συναρτήσεις κατάταξης και αποκατάταξης (\eng{rank}/\eng{unrank}). Για την παραγωγή μικρών πινάκων JSON χρησιμοποιούνται πιο στοχευμένες προβολές, επειδή οι πλήρεις χώροι μεταθέσεων είναι μεγαλύτεροι. Η δεύτερη φάση του αλγορίθμου δύο φάσεων προσθέτει τρεις προβολές: τη μετάθεση γωνιών μεγέθους \(8!\), τη μετάθεση των οκτώ ακμών των εδρών \(U\)/\(D\) μεγέθους \(8!\) και τη μετάθεση των τεσσάρων ακμών της μεσαίας φέτας μεγέθους \(4!\). Οι πίνακες αυτοί παράγονται μόνο με κινήσεις που διατηρούν την ενδιάμεση ομάδα. Για τον επιλυτή Korf, που βασίζεται στην επαναληπτικά εμβαθυνόμενη αναζήτηση A* (\eng{iterative-deepening A*}, IDA*), προστίθεται ξεχωριστή συνδυασμένη γωνιακή συντεταγμένη μεγέθους \(8!\cdot3^7\), η οποία ενώνει μετάθεση και προσανατολισμό γωνιών και αποθηκεύεται σε δυαδικό πίνακα προτύπων (\eng{pattern database}, PDB), κατά το πρότυπο του Korf~\cite{korf1997pattern}.

Η βασική ιδέα είναι ότι μια πραγματική λύση του κύβου προβάλλεται σε λύση κάθε συντεταγμένης· επομένως η απόσταση μιας προβολής από τη λυμένη της τιμή αποτελεί κάτω φράγμα για την απόσταση της πλήρους κατάστασης, σύμφωνα με το επιχείρημα παραδεκτότητας που αναλύθηκε στο Κεφάλαιο~\ref{chap:background}.

\section{Πίνακες κινήσεων}
\label{sec:model-move-tables}
Ο πίνακας κινήσεων μιας συντεταγμένης είναι ένας δισδιάστατος πίνακας. Η γραμμή αντιστοιχεί στην τρέχουσα τιμή της συντεταγμένης και η στήλη σε κίνηση· η τιμή είναι η νέα συντεταγμένη μετά την εφαρμογή της κίνησης. Για παράδειγμα, ο πίνακας CO έχει 2187 γραμμές και 18 στήλες. Η κατασκευή του πίνακα γίνεται με αποκωδικοποίηση της συντεταγμένης σε αντιπροσωπευτική κατάσταση κυβιδίων, εφαρμογή της κίνησης στο μοντέλο κυβιδίων και επανακωδικοποίηση της νέας κατάστασης.

Η διαδικασία αυτή είναι πιο αργή από χειροποίητους μαθηματικούς τύπους μετάβασης, αλλά έχει σημαντικό μεθοδολογικό πλεονέκτημα: συνδέει άμεσα τον πίνακα με το ήδη δοκιμασμένο μοντέλο κινήσεων. Οι δοκιμές συγκρίνουν τις τιμές του πίνακα με την άμεση εφαρμογή κινήσεων σε καταστάσεις κυβιδίων. Έτσι ο πίνακας δεν είναι μαύρο κουτί· προκύπτει από το μοντέλο και μπορεί να αναπαραχθεί.

Οι πίνακες κινήσεων αποθηκεύονται σε JSON μαζί με το όνομα της συντεταγμένης, το είδος του πίνακα, το προφίλ εκτέλεσης, τον σπόρο (\eng{seed}) και τη λίστα κινήσεων. Αν και ο σπόρος δεν αλλάζει το περιεχόμενο ενός αιτιοκρατικού πίνακα κινήσεων, καταγράφεται ώστε η ιχνηλασιμότητα να είναι ενιαία με αυτήν των μετρήσεων. Τα αθροίσματα ελέγχου (\eng{checksums}) στο μανιφέστο (\eng{manifest}) επιτρέπουν στον αναγνώστη να διαπιστώσει αν τα αρχεία έχουν αλλάξει μετά την παραγωγή τους.

\section{Πίνακες αποκοπής}
\label{sec:model-pruning-tables}
Ο πίνακας αποκοπής (\eng{pruning table}) αποθηκεύει για κάθε τιμή συντεταγμένης την ελάχιστη απόσταση της προβολής από τη λυμένη προβολή, με βάση το αντίστοιχο σύνολο κινήσεων. Παράγεται με αναζήτηση κατά πλάτος (\eng{breadth-first search}, BFS) στον γράφο συντεταγμένων, με αφετηρία τη λυμένη συντεταγμένη. Επειδή ο γράφος αυτός είναι πολύ μικρότερος από τον πλήρη χώρο καταστάσεων, η BFS είναι εφικτή για τις προβολές που χρησιμοποιεί η εργασία.

Οι πίνακες αποκοπής χρησιμοποιούνται με δύο τρόπους. Πρώτον, δίνουν παραδεκτά ευρετικά κάτω φράγματα στην IDA*. Η ευρετική συνάρτηση λαμβάνει τη μέγιστη τιμή από τις διαθέσιμες προβολές, τη γωνιακή PDB, τις ακμικές PDB και, όπου χρειάζεται, το κάτω φράγμα λανθασμένων κυβιδίων. Η μέγιστη τιμή παραμένει παραδεκτή, επειδή καμία από τις επιμέρους προβολές δεν υπερεκτιμά. Δεύτερον, στις φάσεις του αλγορίθμου δύο φάσεων οι πίνακες χρησιμοποιούνται ως φασικά κάτω φράγματα για συγκεκριμένους στόχους, όπως ο μηδενικός προσανατολισμός ακμών ή η επίτευξη της ενδιάμεσης ομάδας. Η τετραφασική μέθοδος Thistlethwaite χρησιμοποιεί πίνακες αποστάσεων αντίστοιχου τύπου, αλλά ακριβείς, ανά φάση. Αυτοί δεν λειτουργούν ως ευρετικά κάτω φράγματα αναζήτησης, αλλά κατευθύνουν απευθείας τη μετάβαση κάθε σταδίου προς τον στόχο του.

Οι πίνακες της εργασίας έχουν διαφορετικές κλίμακες και θεμελιώνουν αντίστοιχα διαφορετικούς ισχυρισμούς. Ο πίνακας αποκοπής CO είναι πλήρης για την προβολή CO, όχι για τον κύβο. Η γωνιακή PDB είναι πλήρης για τη συνδυασμένη γωνιακή προβολή του 3×3×3, ενώ κάθε ακμική PDB είναι πλήρης για έξι επιλεγμένες ακμές. Η εργασία επομένως αναφέρεται σε πίνακες αποκοπής προβολών, σε δικούς της πλήρεις γωνιακούς και ακμικούς πίνακες προτύπων και σε παραδεκτά κάτω φράγματα. Δεν αναφέρεται σε πλήρη πίνακα βέλτιστων αποστάσεων, δηλαδή μαντείο (\eng{oracle}), για τον 3×3×3.

\section{Ιχνηλασιμότητα συντεταγμένων}
\label{sec:model-traceability}
Η ιχνηλασιμότητα των συντεταγμένων είναι μέρος της πειραματικής μεθοδολογίας. Κάθε παραγόμενος πίνακας έχει γραμμή στο μανιφέστο με το μέγεθος του πεδίου ορισμού, το πλήθος εγγραφών, τις κινήσεις, το μέγεθος αρχείου, τον χρόνο παραγωγής και το άθροισμα ελέγχου. Οι πίνακες με τα μεταδεδομένα που παρουσιάζονται στο Κεφάλαιο~\ref{chap:implementation} παράγονται από το ίδιο αρχείο. Με αυτόν τον τρόπο δεν υπάρχουν χειροκίνητα αντιγραμμένες τιμές στο κείμενο.

Το ίδιο σκεπτικό ακολουθούν και οι δοκιμές. Οι δοκιμές αμφίδρομης κωδικοποίησης συντεταγμένων ελέγχουν ότι η κωδικοποίηση και η αποκωδικοποίηση συμφωνούν. Οι δοκιμές πινάκων κινήσεων ελέγχουν ότι η μετάβαση μέσω πίνακα συμφωνεί με την απευθείας εφαρμογή κίνησης. Οι δοκιμές πινάκων αποκοπής ελέγχουν ότι η λυμένη συντεταγμένη έχει απόσταση μηδέν και ότι οι αποστάσεις χρησιμοποιούνται ως μη αρνητικά κάτω φράγματα. Οι έλεγχοι αυτοί δεν αποδεικνύουν ότι ολόκληρος ο επιλυτής του 3×3×3 είναι πλήρης, αποδεικνύουν όμως ότι η υποδομή των πινάκων δεν είναι απλή θεωρητική περιγραφή.

Η σχεδίαση αυτή καθιστά σαφές ποια στοιχεία είναι δεδομένα εισόδου και ποια παραγόμενα. Οι βιβλιογραφικές πηγές είναι τεκμηριωμένες εξωτερικές πηγές, ενώ οι πίνακες συντεταγμένων, οι γραμμές των μετρήσεων και οι πίνακες του κειμένου είναι παραγόμενα αρχεία. Η διάκριση αυτή αποτρέπει τη σύγχυση ανάμεσα σε βιβλιογραφικούς αριθμούς, τοπικά αποτελέσματα και χειροκίνητες παρατηρήσεις.
